\chapter{CR Manifolds}
\textcolor{blue}{Note that the wedge $\wedge $  is taken by Alt convention, not Det convention as usual.}

\section{CR manifolds}
Let $M$ be a $m$-dimensional smooth manifold.  Let $n\in \nn$ be an integer s.t. $1\leq n \leq [m/2]$.  
\begin{definition}
    A complex subbundle $T_{1,0}(M)$ of the complexified tangent bundle $T(M) \otimes \mathbb{C}$, of complex rank $n$, is called \textbf{an almost complex CR structure} if 
    $$
T_{1,0}(M) \cap T_{0,1}(M)=(0).
$$ 
Here $ T_{0,1}(M)=\overline{T_{1,0}(M)}$. The integers $n$ and $k=m-2n$ are respectively the CR dimension and CR codimension of the almost CR structure and $(n,k)$ is its type.    
\end{definition}
\begin{remark}[cf. \cite{dragomirDifferentialGeometryAnalysis2006}]
An almost CR manifold of type $(n,0) $ is an almost complex manifold.  
\end{remark}
\begin{definition}
    An almost CR structure is \textbf{formally integrable} if
    $$
\left[\Gamma^{\infty}\left(U, T_{1,0}(M)\right), \Gamma^{\infty}\left(U, T_{1,0}(M)\right)\right] \subseteq \Gamma^{\infty}\left(U, T_{1,0}(M)\right).
$$
If this is the case, $T_{1,0}(M)$ is referred to as a\textbf{ CR structure} and $(M,T_{1,0}(M))$ is a \textbf{CR manifold. } 
\end{definition}
\begin{definition}
    A smooth map $f$ between two CR manifolds is called a \textbf{CR map} if   
    $$
\left(d_{x} f\right) T_{1,0}(M)_{x} \subseteq T_{1,0}(N)_{f(x)},
$$
where $df$ is the $\cc$-linear extension of the differential map.  
\end{definition}
\begin{remark}
    CR manifolds are the objects of a category whose arrows are smooth maps preserving CR structures. The  complex manifolds and holomorphic maps form a subcategory of the category of CR manifolds and CR maps.
\end{remark}

Let $\crm$ be a CR manifold of type $(n,k)$.
\begin{definition}
    The\textbf{ maximal complex} or \textbf{Levi distribution} of $\crm$ is  
    $$
H(M)=\operatorname{Re}\left\{T_{1,0}(M) \oplus T_{0,1}(M)\right\}.
$$
It carries the \textbf{complex structure} $J_{b}: H(M) \rightarrow H(M)$ given by
$$
J_{b}(V+\bar{V})=i(V-\bar{V}).
$$
Formally integrable condition is equivalent to
$$
\begin{gathered}
{\left[J_{b} X, Y\right]+\left[X, J_{b} Y\right] \in \Gamma^{\infty}(U, H(M)),} \\
{\left[J_{b} X, J_{b} Y\right]-[X, Y]=J_{b}\left\{\left[J_{b} X, Y\right]+\left[X, J_{b} Y\right]\right\}}
\end{gathered}
$$
CR map is equivalent to say
$$
\left(d_{x} f\right) H(M)_{x} \subseteq H(N)_{f(x)}
$$
and
$$
\left(d_{x} f\right) \circ J_{b, x}=J_{b, f(x)}^{N} \circ\left(d_{x} f\right).
$$
\end{definition}
  
\begin{definition}
   $ f : M \to N$ is a\textbf{ CR isomorphism} (or a CR equivalence) if $f$ is a smooth  diffeomorphism and a CR map.
\end{definition}

CR manifolds arise mainly as real submanifolds of complex manifolds. Let $V$ be a complex manifold, of complex dimension $N$ , and let $M \subset V$ be a real $m$-dimensional submanifold. Let us set
$$ T_{1,0}(M):=T^{1,0}(V) \cap T(M) \otimes \cc. $$
\begin{proposition}
    If $M$ is a real hypersurface $(m = 2N - 1)$, then$ T_{1,0}(M)$ is a CR structure of type $(N - 1, 1)$. 
\end{proposition}
In general the complex dimension of $T_{1,0}(M)$ defined above may not be constant. Nevertheless, if $ \dim _{\cc}T_{1,0}(M)\equiv n $, then $\crm$  is a CR manifold of type $(n,k)$.
\begin{definition}
    If the CR codimension of $\crm$ and the codimension of $M$ as a real submanifold of $V$ coincide, i.e., $k=2N-m$,  then $\crm$  is termed \textbf{ generic}.
\end{definition}
   

\subsection{Hypersurface type}
 Let 
$$ E_x=\{\omega \in T^{*}_x(M):H(M)_x \subset \ker (\omega )\} $$ 
Assume $M$ is orientable, then the bundle $E$ has a globally defined nowhere vanishing sections $\theta \in \Gamma (E)$.  
Given any \textbf{pseudo-Hermitian structure} $\theta$ on $M$ the \textbf{Levi form} $L_{\theta}$ is defined by
$$
\boxed{L_{\theta}(Z, \bar{W}):=-i(d \theta)(Z, \bar{W})=G_{\theta }(Z, \overline{W}).}
$$

\begin{remark}
The definition of Levi form of type $(n,k)$ could be possible, cf. \cite{dragomirDifferentialGeometryAnalysis2006}. 
\end{remark}

Define a bilinear form (also called \textbf{Levi form})
$$
\boxed{G_{\theta}(X, Y)=(d \theta)\left(X, J_{b} Y\right),}
$$
then
$$
\begin{aligned}
G_{\theta}\left(J_{b} X, J_{b} Y\right)-G_{\theta}(X, Y) & =-(d \theta)\left(J_{b} X, Y\right)-(d \theta)\left(X, J_{b} Y\right) \\
& =\frac{1}{2}\left\{\theta\left(\left[J_{b} X, Y\right]\right)+\theta\left(\left[X, J_{b} Y\right]\right)\right\} \\
& =\frac{1}{2} \theta\left(J_{b}\left\{[X, Y]-\left[J_{b} X, J_{b} Y\right]\right\}\right)=0.
\end{aligned}
$$
Define the $2$-form $\Omega$ by 
$$\boxed{\Omega(X,Y):=g_{\theta}(X,JY)=-d \theta (X,Y).} $$

\begin{definition}
    We say that $\left(M, T_{1,0}(M)\right)$ is \textbf{nondegenerate} if the Levi form $L_{\theta}$ is, the pair $(M, \theta)$ is referred to as a \textbf{pseudo-Hermitian manifold},
    and \textbf{strictly pseudoconvex} if $L_{\theta}$ is positive.
\end{definition}

\begin{remark}
Any two pseudo-Hermitian structure $\theta ,\hat{\theta }$ are related by 
\begin{equation}\label{eq:crtran}
    \hat{\theta }=\lambda \theta 
\end{equation}   
for some non-zero function $\lambda .$ Then Levi form changes according to 
$$ L_{\hat{\theta }}=\lambda L_{\theta }. $$ 
Hence non-degeneracy is a \textbf{CR-invariant} property, i.e., it is invariant under a transformation $\eqref{eq:crtran}.$ 
\end{remark}

\begin{definition}
    We say that a CR map $f: M \rightarrow N$ is a \textbf{pseudo-Hermitian map} if for some $c\in \mathbb{R}$,
    $$f^{*}\theta=c\theta. $$
    If $c=1$, then $f$ is referred to as an \textbf{isopseudo-Hermitian map}.
\end{definition}

\begin{definition}[Webster metric]
    $$
g_{\theta}(X, Y)=G_{\theta}(X, Y), \quad g_{\theta}(X, T)=0, \quad g_{\theta}(T, T)=1
$$
$$
\boxed{g_{\theta}=\pi_{H} G_{\theta}+\theta \otimes \theta.}
$$
\end{definition}
\begin{proposition}[Reeb vector field]
  There is a unique globally defined nowhere zero tangent vector field $T$ on $M$ such that
  $$ \theta (T)=1, \quad T \rfloor \, d \theta =0.
  $$
\end{proposition}



\subsubsection{CR geometry and contact Riemannian geometry}
Let $M$ be a $(2n+1)$-dimensional smooth manifold. Let $(\phi,\xi,\eta)$
be a synthetic object, consisting of a $(1, 1)$-tensor field $\phi: T (M) \to T (M)$, a tangent
vector field $\xi \in \mathcal{X}(M)$, and a differential $1$-form $\eta$ on $M$. Then $(\phi,\xi,\eta)$ is an \textbf{almost contact structure} if
$$
\phi^{2}=-I+\eta \otimes \xi, \quad \phi \xi=0, \quad \eta(\xi)=1, \quad \eta \circ \phi=0 .
$$
An almost contact structure $(\phi,\xi,\eta)$ is said to be \textbf{normal} if (recall $ \eqref{eq:Jtorsion} $)
$$
[\phi, \phi]+2(d \eta) \otimes \xi=0.
$$
A Riemannian metric $g$ on $M$ is said to be \textbf{compatible} with the almost contact structure $(\phi,\xi,\eta)$ if
$$
g(\phi X, \phi Y)=g(X, Y)-\eta(X) \eta(Y).
$$
A synthetic object $(\phi,\xi,\eta,g)$ consisting of an almost contact
structure $(\phi,\xi,\eta)$ and a compatible Riemannian metric $g$ is said to be an \textbf{almost contact metric structure}. 

Given an almost contact metric structure $(\phi,\xi,\eta,g)$ one defines
a $2$-form $\Omega$ by setting $$\Omega(X, Y ) = g(X, \phi Y ).$$  

$(\phi,\xi,\eta,g)$ is said to satisfy the \textbf{contact condition }if $\Omega = d\eta$, and if this is the case, $(\phi,\xi,\eta,g)$ is called a \textbf{contact metric structure }on $M$. 

A contact metric structure $(\phi,\xi,\eta,g)$ that is also normal is called a \textbf{Sasakian structure} (and $g$ is a \textbf{Sasakian metric}).

\begin{example}
    Strictly pseudoconvex CR manifolds are contact Riemannian manifolds, in a natural way.
    In fact,
    $$
\begin{aligned}
J_{b}^{2} & =-I+\theta \otimes T, \\
J_{b} T & =0, \quad \theta \circ J_{b}=0, \quad g_{\theta}(X, T)=\theta(X), \\
g_{\theta}\left(J_{b} X, J_{b} Y\right) & =g_{\theta}(X, Y)-\theta(X) \theta(Y),
\end{aligned}
$$
and
$$\Omega=-d\theta.$$
\end{example}

\begin{theorem}[\cite{dragomirDifferentialGeometryAnalysis2006}]
    Let $(M, T_{1,0} (M), \theta)$ be a pseudo-Hermitian manifold. Then the Webster metric $g_{\theta}$ is Sasakian if and only if the Tanaka-Webster connection $\nabla$ has vanishing \textbf{pseudo-Hermitian torsion}, that is, $\tau=0.$ 
\end{theorem}
\begin{remark}
    The pseudo-Hermitian torsion of $\nabla $,  $\tau$ is defined by 
    $$\boxed{\tau(X)=T_{\nabla}(T,X).}$$ 
    $T_{\nabla}$ is torsion tensor of the Tanaka–Webster connection.
\end{remark}

\begin{theorem}[Tanaka-Webster connection on pseudo-Hermitian manifold]
    There exists a connection $\nabla$ with the following properties:\\
    1. $H(M)$ is parallel w.r.t. $\nabla$;\\
    2. $\nabla g_{\theta}=0,\nabla J=0, \nabla \theta =0,\nabla T=0.$\\
    3. $T_{\nabla}(X,Y)=-2\Omega(X,Y)T, T_{\nabla}(T,JX)=-JT_{\nabla}(T,X)$.
\end{theorem}

\begin{remark}
    We say $T_{\nabla}$ is \textbf{pure} if
    \begin{align}
    T_{\nabla}(Z, W) & =0 \\
    T_{\nabla}(Z, \bar{W}) & =2 i L_{\theta}(Z, \bar{W}) T \label{eq:pure2}\\
    \tau \circ J+J \circ \tau & =0
    \end{align}
\end{remark}

\begin{theorem}[Tanaka-Webster connection]
    Let $(M, T_{1,0}(M),\theta) $ be a nondegenerate CR manifold. Then there is a unique linear connection $ \nabla  $ on $M$ satisfying the following axioms:
    \begin{enumerate}[(i)]
        \item $H(M)$ is parallel w.r.t. $\nabla$.
        \item $ \nabla J =0,\quad \nabla g_{\theta }=0. $
        \item $ T_{\nabla } $ is pure.
    \end{enumerate}
\end{theorem}
\begin{proof}
    Let $ \pi _{\pm}: T_{\cc}M \to T_{1,0}(M)~ (T_{0,1}(M)) $ be the natural projections. Then $ \pi _{-}Z=\overline{\pi _{+} \bar{Z} }$. Note also that 
    $$ T_{\cc}M=T_{1,0}(M) \oplus T_{1,0}(M) \oplus \cc T.$$

    \textbf{Uniqueness}: From $ \eqref{eq:pure2} $, 
    $$ [\bar{Y}, Z]=\nabla_{\bar{Y}} Z-\nabla_Z \bar{Y}+2 i L_\theta(Z, \bar{Y}) T, \quad \forall Y,Z \in T_{1,0}(M) .$$

\end{proof}

\begin{proposition}\label{pr:connection}
    The following formulas hold:
   $$ \boxed{T_{\nabla }=2(\theta  \wedge \tau -\Omega \otimes  T) }$$ 
   $$ \boxed{\nabla ^{\theta }=\nabla +(\Omega -A)T+\tau  \otimes \theta +2 \theta \odot  J }$$
   $$ \trace_{g_{\theta }}(\nabla ^{\theta }-\nabla )=\trace_{G_{\theta }}(\nabla ^{\theta }-\nabla )=0. $$
\end{proposition}


\begin{theorem}
  The curvatures of  TW connection and LC connection  are related by 
  $$ \begin{aligned}
    R^\theta(X, Y) Z=R(X, Y) & Z+(L X \wedge L Y) Z-2 \Omega(X, Y) J Z \\
    & -g_\theta(S(X, Y), Z) T+\theta(Z) S(X, Y) \\
    & -2 g_\theta((\theta \wedge \mathcal{O})(X, Y), Z) T+2 \theta(Z)(\theta \wedge \mathcal{O})(X, Y),
    \end{aligned} $$
    for any $X,Y,Z \in TM$. 
    where 
    $$ S(X,Y)=(\nabla _{X}\tau )Y-(\nabla _{Y}\tau )X ,$$
    $$ \mathcal{O}=\tau ^2+2J \tau -I ,\quad L=\tau +J.$$
\end{theorem}
\begin{remark}
   Note that 
   $$ J^2Y=-Y+\theta (Y)T, \quad \forall Y \in TM. $$ 
\end{remark}

\begin{theorem}[Structure equations for TW connection]
    \begin{equation}
    \begin{aligned}
d\theta & =2\sqrt{-1}h_{\alpha\bar{\beta}}\theta^{\alpha}\wedge\theta^{\bar\beta},  \\
d\theta^{\alpha}& =\theta^\beta\wedge\omega_\beta^\alpha+\theta\wedge\tau^\alpha, & dh_{\alpha\bar{\beta}}=\omega_{\alpha}^{\gamma}h_{\gamma\bar{\beta}}+h_{\alpha\bar{\gamma}}\omega_{\bar{\beta}}^{\bar{\gamma}},  \\
d\omega_\beta^{\alpha}& =-\omega_\gamma^\alpha\wedge\omega_\beta^\gamma+\Pi_\beta^\alpha, 
\end{aligned}
\end{equation}
where
\begin{equation}
\begin{aligned}
\Pi_\beta^\alpha&=R_{\beta\gamma\bar\delta}^\alpha\theta^\gamma\wedge\theta^{\bar\delta}+W_{\beta\gamma}^\alpha\theta^\gamma\wedge\theta-W_{\beta\bar\gamma}^\alpha\theta^{\bar\gamma}\wedge\theta\\&
+2\sqrt{-1}\theta_\beta\wedge\tau^\alpha-2\sqrt{-1}\tau_\beta\wedge\theta^\alpha,
\end{aligned}
\end{equation}
and
\begin{equation}
    W_{\alpha \bar{\gamma}}^{\beta}=h^{\bar{\delta} \beta} A_{\bar{\gamma} \bar{\delta}, \alpha}, \quad W_{\alpha \gamma}^{\beta}=h^{\bar{\delta} \beta} A_{\alpha \gamma, \bar{\delta}}, \quad \tau_{\alpha}=h_{\alpha \bar{\beta}} \tau^{\bar{\beta}}, \quad \theta_{\alpha}=h_{\alpha \bar{\beta}} \theta^{\bar{\beta}}.
\end{equation}
\end{theorem}
\begin{remark}
    \begin{equation}
d\theta^{\bar\alpha}=\theta^{\bar\beta}\wedge\omega_{\bar\beta}^{\bar\alpha}+\theta\wedge\tau^{\bar\alpha}
    \end{equation}
\end{remark}

\section{Divergence formulas}
\begin{definition}
    Define the divergence of a vector $X \in TM$  by 
    $$ \di (X) \Psi  = \mathcal{L}_{X} \Psi ,$$
    where $\Psi $ is a volume form. 
\end{definition}
\begin{remark}
   We can take 
   $$ \Psi = \theta \wedge (d \theta )^n .$$ 
\end{remark}
\begin{remark}
    By 3rd relation of Proposition \ref{pr:connection}, 
   $$ \di (X)= \trace_{g_{\theta }}\{Y \in TM \to \nabla _{Y}X\} .$$ 
   Note that 
   $$ \begin{aligned}
   g_{\theta }(\nabla ^{\theta }_{e_{A}}X,e_{A})= g_{\theta }(\nabla ^{\theta }_{T_{\alpha }}X,T_{\bar \alpha })+g_{\theta }(\nabla ^{\theta }_{T_{\bar\alpha }}X,T_{ \alpha })+g_{\theta }(\nabla ^{\theta }_{T}X,T),
   \end{aligned}  $$
   assuming $\{T_{\alpha }\} $ is orthonormal and $T_{\alpha }=\frac{1}{\sqrt{2}}(e_{\alpha }-iJ e_{\alpha }) $.  
\end{remark}
\begin{exercise}
    For any complex vector $Z \in T_{\cc}(M)$, 
    $$ \overline{\di (Z)}=\di (\bar{Z}) .$$
\end{exercise}
\begin{definition}
    Define  the divergence of a 1-form $\sigma $ by 
    $$ \di (\sigma )=\di (X_{\sigma }), $$ 
    where $X_{\sigma }$ is the dual vector of $\sigma $ w.r.t. $g_{\theta }$, i.e., $g_{\theta }(X_{\sigma },Y)=\sigma (Y), \forall  Y \in TM $.    
\end{definition}
\begin{exercise}
$$ \di (Z)=T_{\alpha }(Z^{\alpha })+Z^{\beta } \omega _{\beta }^{\alpha }(T_{\alpha }), \quad \forall  Z=Z^{\alpha }T_{\alpha }\in T_{1,0}(M).$$
$$ \di (\sigma )=\sigma _{\alpha ,\bar{\beta }}h^{\alpha \bar{\beta }}. $$
\end{exercise}

\begin{definition}
    Sublaplacian for functions is defined by 
    $$ \Delta _b(u):=-\di(\nabla^{\mathcal{H}} u)=-(u_{\alpha \bar{\beta }}h^{\alpha \bar{\beta }}+u_{\bar{\alpha }\beta }h^{\bar{\alpha }\beta }) ,$$
    where $\nabla ^{\mathcal{H}}u$ is the horizontal component of the gradient $\nabla u$.  
\end{definition}
\begin{exercise}
    The usual Laplacian is given by 
    $$ \Delta  u= -(u_{\alpha \bar{\beta }}h^{\alpha \bar{\beta }}+u_{\bar{\alpha }\beta }h^{\bar{\alpha }\beta }+u_{00}). $$
\end{exercise}










\chapter{Geometry of CR-Submanifolds}

\section{CR-submanifolds and CR-structures}
Let $N$ be a $n$-dimensional \textit{almost Hermitian manifold} with almost complex structure $J$ and with Hermitian metric $g$. Let $M$ be a \textit{real} $m$-dimensional Riemannian manifold isometrically immersed in $N$.

\begin{definition}
    The differential geometry of $M$ depends on the behaviour of the tangent bundle of $M$ relative to the action of the almost complex structure $J$. Thus $M$ is called a \textbf{\textit{complex} (holomorphic)} submanifold if $T_xM$ is invariant by $J$, i.e., we have
$$
J(T_xM) = T_xM, \quad \text{for each } x \in M.
$$
Also, we say that $M$ is a \textbf{\textit{totally real} (anti-invariant)} submanifold of $N$ if we have
$$
J(T_xM) \subseteq T_xM^{\perp}, \quad \text{for each } x \in M.
$$
These two classes of submanifolds have been extensively investigated in the last decade from different viewpoints.
\end{definition}

\begin{definition}
    In 1978 ``we" (refer to the author of \cite{bejancuGeometryCRSubmanifolds1986}) initiated a study of the differential geometry of a new class of submanifolds situated between the above two classes, called CR-submanifolds. More precisely, $M$ is said to be a \textbf{CR-submanifold} of $N$ if there exists a differentiable distribution
\[
D: x \mapsto D_x \subseteq T_xM,
\]
on $M$ satisfying the following conditions:
\begin{enumerate}
    \item[(i)] $D$ is holomorphic, i.e., $J(D_x) = D_x$ for each $x \in M$,
    \item[(ii)] the complementary orthogonal distribution
    \[
    D^{\perp}: x \mapsto D^{\perp}_x \subseteq T_xM^{\perp},
    \]
    is anti-invariant, i.e., $J(D^{\perp}_x) \subseteq (T_xM)^{\perp}$ for each $x \in M$.
\end{enumerate}
We denote by $p$ the complex dimension of the distribution.
D and by $q$ the real dimension of the distribution $D^\perp$. Then for $q = 0$ (resp. $p = 0$) a CR-submanifold becomes a \textbf{complex submanifold (resp. totally real submanifold)}. If $q = \dim T_xM^\perp$ the CR-submanifold is called an \textbf{anti-holomorphic submanifold}. A \textbf{proper CR-submanifold} is a CR-submanifold which is neither a complex submanifold nor a totally real submanifold.
\end{definition}

\begin{example}
    Each real hypersurface $M$ of $N$ (with $n \geq 2$) is a proper CR-submanifold. In fact we define
\[
D_x^\perp : x \mapsto D_x^\perp = J(T_xM^\perp)
\]
and take $D$ as the complementary orthogonal distribution to $D^\perp$ in $TM$. Thus $M$ is endowed with a pair of distributions $(D, D^\perp)$ satisfying the conditions of the definition of a CR-submanifold. Moreover, we have $\dim _{\mathbb{R}} D_x^\perp = 1$ and $\dim_{\mathbb{C}} D_x = n - 1$.

Hence, $M$ is a proper CR-submanifold.
\end{example}


\begin{remark}
    It is easily seen that on a CR-submanifold the 
distribution \( D \) (resp. \( D^{\perp} \)) is the maximal distribution 
invariant by \( J \) (resp. anti-invariant by \( J \)), i.e., if \( D' \) 
(resp. \( D'' \)) is an invariant (resp. anti-invariant) distribution on \( M \), 
then we have \( D'_{x} \subseteq D_{x} \) (resp. \( D''_{x} \subseteq D_{x} \)), for each \( x \in M \).
\end{remark}

\begin{definition}
    Now let \( M \) be an arbitrary Riemannian manifold isometrically immersed in an almost Hermitian manifold \( N \). 
For each vector field \( X \) tangent to \( M \) we put
\begin{equation}
    JX = \phi X + \omega X,
\end{equation}
where \( \phi X \) and \( \omega X \) are respectively the tangent part and the normal part of \( JX \). Also, for each vector field \( V \) normal to \( M \) we put
\begin{equation}
    JV = BV + CV,
\end{equation}
where \( BV \) and \( CV \) are respectively the tangent part and the normal part of \( JV \).
\end{definition}


\begin{theorem}
     The submanifold \( M \) of \( N \) is a CR-submanifold if and only if we have
\[
\text{rank}(\phi) = \text{constant}, \omega\phi = 0.
\]
or, if and only if
\[
\text{rank}(B) = \text{constant}, \phi B = 0.
\]
\end{theorem}
\begin{remark}
Note that 
\begin{equation}
    \begin{aligned}
        &\phi X=JPX\\
        &\omega X=JQX
    \end{aligned}
\end{equation}
    for any $X\in TM$, $P, Q$ is the projection to $D, D^{\perp}$ respectively.

    Also, 
\begin{equation}
    \begin{aligned}
        &D=Im (\phi)=ker (\omega)\\
        &D^{\perp}=Im (B)=ker (\phi).
    \end{aligned}
\end{equation}
\end{remark}

\begin{remark}
     Now, suppose \( M \) is a CR-submanifold of the almost Hermitian manifold \( N \). 
    \begin{equation}
        \begin{aligned}
            \phi^2 = -P\\
            \phi^3 + \phi = 0\\
            C^3 + C = 0.
        \end{aligned}
    \end{equation}
\end{remark}


\begin{proposition}
    On each CR-submanifold \( M \) the vector bundle 
morphisms \( \phi \) and \( C \) define \( f \)-structures on \( TM \) and \( T^{\perp}M \) 
respectively.
\end{proposition}

\begin{theorem}[Blair-Chen]
    A CR-submanifold of a \textbf{Hermitian} manifold is a CR-manifold.
\end{theorem}

\section{f-structure}

Let \( M_n \) be an \( n \)-dimensional differentiable manifold of class \( C^{\infty} \) and 
let there be given a tensor field \( f \neq 0 \) of type \( (1,1) \) and of class \( C^{\infty} \) such that
\begin{equation}
f^3 + f = 0.
\label{eq:f_cubed_plus_f}
\end{equation}

\begin{theorem}[\cite{YANO1982251}]
    For a tensor field \( f \neq 0 \) satisfying \eqref{eq:f_cubed_plus_f}, the operators
\begin{equation}
l = -f^2, \quad m = f^2 + 1,
\label{eq:operators_l_m}
\end{equation}

\( 1 \) denoting the identity operator, applied to the tangent space at a point of the manifold are complementary projection operators.

Moreover,
\begin{equation}
fl = lf = f,
\label{eq:fl_lf_f}
\end{equation}
\begin{equation}
fm = mf = 0, lm=ml=0,
\label{eq:fm_mf}
\end{equation}
\begin{equation}
f^2 l = -l,
\label{eq:f2l}
\end{equation}
\begin{equation}
f^3 m = 0,
\label{eq:f3m}
\end{equation}
that is, \( f \) acts on \( L \) as an almost complex structure operator and on \( M \) as null operator.
\end{theorem}

\begin{remark}
    Thus if there is given a tensor field \( f \neq 0 \) satisfying $\eqref{eq:f_cubed_plus_f}$, then there exist 
complementary distributions \( L \) and \( M \) corresponding to the projection operators 
\( l \) and \( m \) respectively. 

Moreover, one can show that if the rank of \( f \) is constant and is equal to \( r \), 
then the dimensions of \( L \) and \( M \) are \( r \) and \( n-r \) respectively. We call such a structure 
an \( f \)-structure of rank \( r \).
\end{remark}

Let me mention that
\begin{theorem}[\cite{YANO1982251}]
     A necessary and sufficient condition for an \( n \)-dimensional 
manifold to admit a tensor field \( f \neq 0 \) of type \( (1,1) \) and of rank \( r \) such 
that \( f^3 + f = 0 \) is that \( r \) is even: \( r = 2m \) and the group of tangent bundle 
of the manifold be reduced to the group \( U(m) \times O(n-2m) \).
\end{theorem}


\section{Integrability of Distributions on a CR-Submanifold}
First of all, we have 
\begin{equation}
    [\phi, \phi](X, Y) = [\phi X, \phi Y] + \phi^2[X, Y] - \phi([\phi X, Y]) - \phi([X, \phi Y]).
\end{equation}
and
\begin{equation}
    [J, J](X, Y) = [\phi, \phi](X, Y) - \Omega([\phi X, Y] + [X, \phi Y]).
\end{equation}

\begin{theorem}[Bejancu]
Let \( M \) be a CR-submanifold of an 
almost Hermitian manifold \( N \). Then the distribution \( D \) is 
integrable if and only if
\begin{equation}
[J, J](X, Y)^{\perp} = 0
\end{equation}
and
\begin{equation}
Q(\phi, \phi)(X, Y) = 0, \text{ for any } X, Y \in \Gamma(D).
\end{equation}
\end{theorem}

\begin{corollary}
    Let $M$ be a CR-submanifold of a \textbf{Hermitian} manifold $N$. The distribution $D$ is integrable if and only if the Nijenhuis tensor of $\phi$ vanishes identically on $D$.
\end{corollary}

\begin{theorem}
    Let $M$ be a CR-submanifold of an almost Hermitian manifold $N$. The distribution $D$ is integrable if and only if the Nijenhuis tensor of $\phi$ vanishes identically on $D^{\perp}$.
\end{theorem}

\section{CR-submanifold of a nearly Kaehlerian manifold}
....

\section{\texorpdfstring{$\phi$}{phi}-Connections on a CR-Submanifold and CR-Products of Almost Hermitian Manifolds} 

\begin{definition}
    A linear connection $\nabla$ on $M$ is called a \textbf{$\phi$-connection} if $\phi$ is covariantly constant with respect to this connection, that is, we have
    \begin{equation}
        \nabla_X \phi =0, \forall X\in TM
    \end{equation}
\end{definition}

\begin{theorem}
    Both distributions $D$ and $D^{\perp}$ on a CR submanifold of an almost Hermitian manifold are parallel with respect to any $\phi$-connection.
\end{theorem}

\begin{definition}
    We say that a CR-submanifold \( M \) of an almost Hermitian manifold \( N \) 
is a \textbf{CR-product} if both distributions \( D \) and \( D^{\perp} \) are integrable 
and \( M \) is locally a Riemannian product \( M_1 \times M_2 \), where \( M_1 \) 
is a leaf of \( D \) and \( M_2 \) is a leaf of \( D^{\perp} \). 
A CR-product with \( D \neq \{0\} \) and \( D^{\perp} \neq \{0\} \) is called a 
\textbf{proper CR-product.}
\end{definition}

\begin{example}[An counterexample]
    There exist no proper CR submanifolds of $S^6$ with both distributions $D$ and $D^{\perp}$ parallel with respect to Levi-Civita connection. However, Sekigawa constructed  an example of a 3-dimensional proper CR-submanifold such that both distributions $D$ and $D^{\perp}$ are integrable.
\end{example}

\begin{theorem}
    Let \( M \) be a CR-submanifold of an 
almost Hermitian manifold \( N \). If the Levi-Civita connection 
on \( M \) is a \( \phi \)-connection, then \( M \) is a CR-product.
\end{theorem}
Recall that
\begin{theorem}
    Both distributions $D$ and $D^{\perp}$ are parallel with respect to Levi-Civita connection  if and only if they are integrable and their leaves are totally geodesic in $N$.
\end{theorem}

\section{CR-submanifolds of Kaehlerian manifolds}

\begin{theorem}
  Let \( M \) be a CR-submanifold of a K\"ahlerian manifold \( N \). Then we have
\begin{enumerate}
    \item[(i)] the distribution \( D^{\perp} \) is integrable,
    \item[(ii)] the distribution \( D \) is integrable if and only if the second fundamental form of \( M \) satisfies
    \[
    h(X, JY) = h(Y, JX), \text{ for any } X, Y \in \Gamma(D).
    \]
\end{enumerate}
\end{theorem}
\begin{remark}
    Blair and Chen have constructed  an example of a CR-submanifold of a Hermitian manifold for which $D^{\perp}$ is not integrable.
\end{remark}


